\section{The Probe on Real Driving Logs}
\label{sec:navsim-probe}

The synthetic study in Section~\ref{sec:synthetic} answered a narrow question,
namely whether a flow-matching temporal transformer can denoise a whole
trajectory at once and reproduce a target distribution, under conditions where
nothing else could go wrong: analytic data, zeroed conditioning, and every
episode the same length. This section describes what the system looks like once
those three simplifications are removed and the same generative core is attached
to the frozen Orbis latent on real NAVSIM driving logs. This is the part the
earlier sections put off, and it is where most of the engineering actually sits.

Three properties of real logs drive most of what follows. The conditioning is no
longer a scalar but a \emph{spatial grid of tokenizer latents}, one grid per
observed frame, and it has to be reduced to something a trajectory transformer
can take in without swamping the two-channel velocity signal. Episodes have
\emph{different lengths}, since a NAVSIM scene is only as long as the run of
frames whose camera images happen to be on disk, so batches are ragged. And the
behaviour we care about, turning or branching at an intersection, is rare, so the
pipeline has to keep whole episodes instead of short fixed windows, and the
evaluation has to be able to pick out the episodes where a distribution over
futures means anything at all. Figure~\ref{fig:navsim-pipeline} shows the
pipeline end to end.

\begin{figure}[htbp]
\centering
\includegraphics[width=\textwidth,height=0.38\textheight,keepaspectratio]{../diagrams/diagram-navsim-pipeline.pdf}
\caption{The NAVSIM probing pipeline. Everything above the dashed arrow runs
once: the driving logs are scanned into an episode index, each episode is decoded
into front-camera frames at $2$\,Hz, and the frozen stage-1 tokenizer turns those
frames into per-frame latents written to one HDF5 file per episode. Everything
below is the training loop. A window dataset reads a whole episode's trajectory
but only the leading image latents, a ring buffer hides the HDF5 read latency, a
padding collate assembles ragged episodes into a rectangular batch with explicit
lengths, and the flow-matching probe is trained and evaluated on that batch with
the padded frames masked out.}
\label{fig:navsim-pipeline}
\end{figure}

\subsection{From driving logs to episodes}
\label{sec:navsim-episodes}

Two source datasets sit behind a single configuration switch. They differ only in
how an \emph{episode} is enumerated; the per-frame dictionary is the same in
both, so all the geometry code is shared. The \textbf{navhard} split
(\texttt{navhard\_two\_stage}) gives one episode per meta pickle, taken as its
longest scene, which is around five frames. That is enough for a sanity check and
far too short for a trajectory distribution. The \textbf{navtrain} split (the
\texttt{trainval} logs plus the \texttt{trainval} sensor blobs) is the one the
experiments use: each log pickle spans many scenes of roughly forty frames at
$2$\,Hz, and one episode is one scene.

The complication with navtrain is that the sensor download is usually
\emph{partial}. The log pickles list every frame, but only some of the
corresponding front-camera JPEGs are actually on disk. Instead of dropping whole
scenes or quietly returning black images, we define an episode as the
\emph{longest contiguous run of frames within a scene whose front-camera image is
present}. This keeps the temporal spacing uniform, which matters because a
velocity is a finite difference and a hole in the frame sequence would silently
double one step. The cost is that some episodes get shorter. Scenes with fewer
than \texttt{min\_frames}$\,=5$ usable frames are thrown away. The scan touches
the whole sensor tree, so the resulting (log, scene) $\to$ present-frames mapping
is cached to \texttt{data/navtrain\_episode\_index.pkl} and only rebuilt when
that file is deleted, for instance after downloading more blobs. The cache used
for the runs here holds $9\,577$ episodes.

Each frame contributes a global ego pose: the translation $(x,y)$ from
\texttt{ego2global\_translation}, and a yaw recovered from the
\texttt{ego2global\_rotation} quaternion $[w,x,y,z]$ as
\begin{equation}
\psi = \operatorname{atan2}\!\big(2(wz+xy),\; 1-2(y^2+z^2)\big).
\label{eq:yaw}
\end{equation}
The whole episode is then expressed in the local frame of its \emph{first} frame,
by subtracting the initial position and rotating by $-\psi_0$,
\begin{equation}
p_t \;=\; R(-\psi_0)\,\big(g_t - g_0\big),
\qquad
R(-\psi_0)=\begin{pmatrix}\cos\psi_0 & \sin\psi_0\\ -\sin\psi_0 & \cos\psi_0\end{pmatrix},
\label{eq:local-frame}
\end{equation}
so every trajectory starts at the origin heading along $+x$. The synthetic
B\'ezier curves satisfied the same convention by construction, which is why the
same normalization and the same plotting code can be reused unchanged.

The model again works in velocity space, with one representational difference
that is easy to misread when comparing the two code paths, so it is worth being
explicit. The synthetic data used a \emph{forward} difference with the last entry
zero-padded, while the NAVSIM loader uses a \emph{backward} difference with the
first entry zeroed,
\begin{equation}
v_0 = 0,\qquad v_t = p_t - p_{t-1}\ \ (t\ge 1).
\label{eq:navsim-velocity}
\end{equation}
The advantage is that the inverse is an unshifted cumulative sum,
$\sum_{i\le t} v_i = p_t - p_0$, so integrating the sampled velocities and
re-anchoring at the first frame recovers the positions exactly with no
off-by-one correction. It also means the first predicted entry is structurally
zero, so the model spends one output slot on a constant. That is harmless, but it
is worth remembering when reading per-frame error curves.

The train/validation split is a deterministic function of the scene token rather
than of the file order. A token is hashed with MD5, reduced modulo $1000$, and
assigned to validation if the bucket comes out below
$1000\cdot\texttt{val\_fraction}$, with $\texttt{val\_fraction}=0.1$. Since the
hash depends only on the token, the split survives re-caching, re-indexing and
any change in the number of downloaded blobs, and the qualitative figures use the
same episodes as the logged validation metrics.

\subsection{Caching the frozen tokenizer latents}
\label{sec:navsim-cache}

Encoding images with the stage-1 tokenizer inside the training loop is possible.
A configuration flag (\texttt{encode\_images}) selects it, and the tokenizer is
then constructed lazily on the correct device while the batch carries raw images.
But it is wasteful, because the encoder is frozen and its output for a given
frame never changes. The default path encodes every episode once, in chunks of
$100$ frames, and writes one HDF5 file per episode holding both tokenizer
streams, \texttt{encoded\_q\_rec} and \texttt{encoded\_q\_sem}, each of shape
$(T, C, H_{\text{lat}}, W_{\text{lat}})$, along with the trajectory, the
velocity, the frame rate and the scene token. With the $\times16$ spatial
compression from Section~\ref{sec:orbis} and the $256\times256$ centre crop the
loader uses, the latent grid is $16\times16$ positions with $C=16$ channels.

The probe consumes the \emph{semantic} stream \texttt{encoded\_q\_sem}. The
semantic encoder is the one regularized towards DINOv2 features, so it is the
stream that should carry scene structure such as lanes, vehicles and free space,
whereas the detail stream is pushed towards the low-level appearance needed for
faithful reconstruction. We cache the reconstruction stream anyway, since it
costs almost nothing extra, so the choice can be ablated later without
re-encoding the dataset.

\subsection{Whole episodes, ragged batches, and the frame mask}
\label{sec:navsim-padding}

The first version of the training loop cut a fixed window of
\texttt{num\_frames} frames out of each episode. That made every batch
rectangular and every tensor shape static, and it also threw away most of the
data. With episodes of about forty frames and a window of eight, the model never
saw a trajectory long enough to contain a complete manoeuvre, and the longest and
most interesting episodes were truncated just like the shortest ones. Setting
\texttt{num\_frames}$\,=0$ switches the dataset to serving \emph{whole}
episodes, at which point the batch becomes ragged.

Ragged batches are handled at the collate level, not inside the model. A custom
collate function pads every time-major tensor of the sample dictionary with zeros
up to the longest episode in the batch, using \texttt{pad\_sequence} so that
trailing dimensions are left alone, and records the true lengths in
\texttt{batch["metadata"]["length"]}. Padding is applied per key independently.
That matters because two different time axes live in one sample: the trajectory
and velocity have length $T$, while the image latents have length
$T_{\text{img}}\le\texttt{context\_images}$, and a very short episode can end up
with fewer image frames than the nominal image context.

The model turns the recorded lengths into a boolean frame mask
$m\in\{0,1\}^{B\times T}$ with $m_{b,t}=1$ iff $t<\ell_b$. The mask is used for
\emph{selection} and not for truncation, which is the part that is easy to get
wrong: the network still predicts all $T_{\max}$ frames of every sample, and the
mask is what removes the padded tail from the loss, from the online velocity
statistics and from every validation metric. Writing the flow-matching residual
for episode $b$ and frame $t$ as $r_{b,t}=\hat u_{b,t}-u_{b,t}\in\mathbb{R}^2$,
the objective is a masked mean over the frames that are both predicted (past the
trajectory context $c$) and valid,
\begin{equation}
\mathcal{L}
= \frac{\displaystyle\sum_{b}\sum_{t\ge c} m_{b,t}\,\tfrac{1}{2}\lVert r_{b,t}\rVert_2^2}
       {\displaystyle\sum_{b}\sum_{t\ge c} m_{b,t}},
\label{eq:masked-loss}
\end{equation}
This is implemented without a per-sample loop, so variable lengths cost one
elementwise multiply. Figure~\ref{fig:navsim-mask} shows the layout of a padded
batch.

\begin{figure}[htbp]
\centering
\includegraphics[width=0.86\textwidth,height=0.34\textheight,keepaspectratio]{../diagrams/diagram-navsim-mask.pdf}
\caption{A padded batch of three episodes. Every episode is stored at the full
padded length $T_{\max}=8$ and the network predicts all $T_{\max}$ velocity steps
for all three. The recorded lengths $\ell_b$ become the frame mask $m_{b,t}$,
which removes the padded tail from the loss of
Equation~\eqref{eq:masked-loss}, from the exponential-moving-average velocity
statistics and from the validation metrics. Only the leading
\texttt{context\_images} frames, here two, carry an image latent, and those are
the only latents read from disk.}
\label{fig:navsim-mask}
\end{figure}

Reading whole episodes would multiply the amount of latent data pulled off disk
if it were not for the second half of the change. The dataset slices the HDF5
arrays so that \emph{only} the leading \texttt{context\_images} latents are read,
while the trajectory, which is two floats per frame, is always read in full. The
denoiser never looks at an image beyond the context, so nothing is lost, and the
I/O per sample becomes essentially independent of episode length. What latency is
left gets absorbed by an in-memory ring buffer, which serves repeated samples and
pulls \texttt{replace\_count}$\,=10$ fresh episodes from the source dataset every
$40$ served samples, growing to at most $20\,000$ cached episodes. The buffer
trades a controlled amount of sample staleness for a GPU that is never starved by
HDF5 decompression.

\subsection{What changed in the image aggregator}
\label{sec:navsim-aggregator}

The conditioning path has been revised more than anything else, and that revision
is the main architectural difference between the model described in
Section~\ref{sec:probe} and the one we actually train.

\paragraph{The original design: one pooled token per frame, concatenated.}
In the first version the aggregator reduced each frame's latent grid to a
\emph{single} token. The $16$ tokenizer channels were lifted to the model width
$D$ by a linear map, the $H_{\text{lat}}\!\times\!W_{\text{lat}}$ positions were
flattened into a sequence with a 2D sinusoidal position embedding, and one
cross-attention query pooled that sequence into a single $D$-dimensional vector,
with the frame embedding entering as the AdaLN conditioning of the pooling block.
The result was a tensor of shape $(B,T,D)$, one image vector per frame, whose
entries beyond the image context were explicitly zeroed. That vector was then
\emph{concatenated to the per-point feature} of every trajectory point, so the
input to the pointwise MLP had width $2+3D$: two velocity channels, the
diffusion-time embedding, the frame embedding and the frame's image token.

Three problems with this only show up on real data. The capacity bottleneck is
severe, since an entire street scene is squeezed into $D=64$ numbers before the
trajectory model sees any of it, so the probe has no way to ask different
questions of different parts of the image. The alignment is rigid: because the
image token is concatenated per frame, image information reaches trajectory point
$t$ only if frame $t$ carried an image, which forces the awkward zeroing of all
non-context frames and burns an input slot on a constant zero vector for most of
the episode. And the interface is brittle, because the width of the first linear
layer depends on how image evidence is supplied, so changing the number of
observed frames changes the parameter shapes.

\paragraph{The current design: several tokens per frame, as attention keys.}
The aggregator now emits $N=\texttt{image\_tokens}$ tokens per observed frame
instead of one. The same cross-attention block is used with $N$ queries rather
than one, giving $(B,T_{\text{img}},N,D)$, and a 1D sinusoidal position embedding
over the $N$ slots is added so the slots stay distinguishable once they are
flattened together with the frames. Only the leading
$T_{\text{img}}=\texttt{context\_images}$ frames get aggregated at all, so the
zeroing of later frames goes away; later frames are simply never encoded. The
resulting tokens are flattened into one sequence of length
$T_{\text{img}}\!\cdot\!N$.

The important part is that these tokens no longer enter the per-point features.
They are prepended to the key/value sequence of every temporal block, so the
per-point input shrinks to width $2+2D$ and the image evidence reaches the
trajectory through \emph{cross-attention} rather than concatenation. Each
predicted frame can now attend to whichever image tokens it needs, the number of
observed frames and the number of tokens per frame can both change without
touching a parameter shape, and the conditioning bandwidth is set by
$T_{\text{img}}\cdot N$ instead of by the model width. In the configuration used
here, $\texttt{context\_images}=2$ and $\texttt{image\_tokens}=9$ give $18$ image
tokens per episode. Figure~\ref{fig:navsim-aggregator} contrasts the two
designs.

There is one asymmetry in the current implementation that should be recorded.
Slot identity enters the image tokens \emph{additively}, through the 1D position
embedding over the $N$ slots, while frame identity enters only
\emph{multiplicatively}, through the AdaLN modulation of the pooling block by the
frame embedding. With $T_{\text{img}}=2$ the two observed frames are distinguished
by a modulation of the pooling rather than by an explicit temporal tag on the
tokens themselves. Adding a frame embedding to the tokens would make the temporal
ordering of the image context explicit, and that is the obvious next refinement
if the image context is ever made longer. One more detail, since
Section~\ref{sec:probe} describes the probe in the abstract: the pooling queries
are \emph{fixed} sinusoidal position embeddings, not learned parameters. What is
learned is the projection of those queries, the key/value projections of the
latent grid, the AdaLN modulation and the feed-forward branch. The query tokens
themselves are deterministic slot identifiers.

\begin{figure}[htbp]
\centering
\begin{subfigure}[t]{0.47\textwidth}
\centering
\includegraphics[width=\linewidth,height=0.32\textheight,keepaspectratio]{../diagrams/diagram-aggregator-old.pdf}
\caption{Before: a single pooled token per frame, concatenated to every
trajectory point.}
\label{fig:navsim-aggregator-old}
\end{subfigure}
\hfill
\begin{subfigure}[t]{0.47\textwidth}
\centering
\includegraphics[width=\linewidth,height=0.32\textheight,keepaspectratio]{../diagrams/diagram-aggregator-new.pdf}
\caption{After: $N$ tokens per observed frame, entering the attention stack as
keys and values.}
\label{fig:navsim-aggregator-new}
\end{subfigure}
\caption{The revision of the image aggregator. Pooling the latent grid into one
vector per frame (left) pushed the whole scene through a $D$-dimensional
bottleneck, tied image evidence to the frame index by concatenation, and required
non-context frames to be zeroed. Emitting $N$ tokens per observed frame and
routing them through cross-attention (right) widens the conditioning bandwidth to
$T_{\text{img}}\cdot N$ tokens, lets every predicted frame query the image context
freely, removes the zero-padding of non-context frames, and makes the per-point
input width independent of how much image evidence is supplied.}
\label{fig:navsim-aggregator}
\end{figure}

\subsection{The denoiser as a whole}
\label{sec:navsim-denoiser}

With the aggregator settled, the rest of the denoiser is a small
sequence-to-sequence transformer whose only trainable inputs are the noisy
velocities and two scalar embeddings. Given the noisy velocity trajectory
$\hat v\in\mathbb{R}^{B\times T\times2}$ at diffusion time $t$, the network forms
a diffusion-time embedding $e_{\text{time}}\in\mathbb{R}^{B\times D}$ broadcast
across frames, and a frame embedding $e_{\text{frame}}\in\mathbb{R}^{T\times D}$
computed from the \emph{normalized} index $\operatorname{arange}(T)/T$ so that it
is well defined for any episode length. Each point is then described by
\begin{equation}
x_i=\big[\underbrace{\hat v_i}_{2},\;\underbrace{e_{\text{time}}}_{D},\;
\underbrace{e_{\text{frame},i}}_{D}\big]\in\mathbb{R}^{2+2D},
\label{eq:navsim-input}
\end{equation}
and split at the trajectory-context boundary: the first
$c=\texttt{context\_traj}$ points go through a \emph{context} MLP and the
remaining $T-c$ through a \emph{prediction} MLP. There are two MLPs rather than
one because the two groups do different jobs. The context tokens are known and
never updated; the prediction tokens are the query stream the stack refines.

Each temporal block is a DiT-style cross-attention block with AdaLN-Zero
conditioning on $e_{\text{time}}$, as in Section~\ref{sec:syn-model}, but here it
is used in a genuinely \emph{cross}-attentive configuration. The query is the
current prediction stream, and the key/value sequence is the concatenation of the
image tokens, the trajectory-context tokens and the prediction tokens,
\begin{equation}
\text{pred}\;\leftarrow\;
\text{TemporalBlock}\big(\underbrace{[\,\text{img}\;\|\;\text{ctx}\;\|\;\text{pred}\,]}_{\text{keys/values}},\;
\underbrace{\text{pred}}_{\text{queries}},\; e_{\text{time}}\big),
\label{eq:navsim-block}
\end{equation}
repeated for every block in the stack. The key sequence is rebuilt from the
updated predictions at each layer, so the image and context tokens stay static
while the prediction stream is refined. That is closer to a decoder stack with a
fixed memory than to plain self-attention. A LayerNorm and a linear map project
the final prediction stream back to two channels, which gives the velocity field
on the predicted frames only, of shape $(B,T-c,2)$.
Figure~\ref{fig:navsim-denoiser} shows the assembled network.
Figure~\ref{fig:navsim-sketch} is a simpler version of the same picture, next
to the tokenizer that produces the image latents.

The configuration used here sets $D=64$ with $8$ attention heads and a
feed-forward width of $2D$, and the stack currently holds two temporal blocks
(reduced from four when the image tokens were added, which lengthened the key
sequence and therefore the cost of each block). With
$\texttt{context\_traj}=0$, $\texttt{context\_images}=2$ and
$\texttt{image\_tokens}=9$, a block attends over $18+T$ keys and produces $T$
queries.

\begin{figure}[htbp]
\centering
\includegraphics[width=\textwidth,height=0.42\textheight,keepaspectratio]{../diagrams/diagram-navsim-denoiser.pdf}
\caption{The trajectory denoiser on real data. The image branch (top) turns the
cached tokenizer latents of the $T_{\text{img}}$ observed frames into
$T_{\text{img}}\cdot N$ tokens. The trajectory branch (bottom left) encodes each
point from the noisy velocity, the diffusion-time embedding and the frame
embedding, splitting at the trajectory-context boundary. Every temporal block
takes the concatenated image, context and prediction tokens as keys and values
and the prediction stream as queries, conditioned on the diffusion time through
AdaLN-Zero, and the refreshed prediction stream re-enters the key sequence at the
next layer. A LayerNorm and a linear head emit the flow-matching velocity field
on the predicted frames only.}
\label{fig:navsim-denoiser}
\end{figure}

\begin{figure}[htbp]
\centering
\begin{subfigure}[t]{0.38\textwidth}
\centering
\includegraphics[width=\linewidth]{../diagrams/diagram-tokenizer.pdf}
\caption{The hybrid tokenizer. The semantic and detail encoders each give an
$18\times32$ grid with $16$ channels, and the two grids are concatenated along
the channels.}
\label{fig:navsim-sketch-tokenizer}
\end{subfigure}
\hfill
\begin{subfigure}[t]{0.58\textwidth}
\centering
\includegraphics[width=\linewidth]{../diagrams/diagram-denoiser-sketch.pdf}
\caption{The denoiser. The spatial block turns each context frame into tokens,
which join the noisy trajectory points in the temporal block. The diffusion
time step comes in through AdaLN-Zero.}
\label{fig:navsim-sketch-denoiser}
\end{subfigure}
\caption{Simplified view of the tokenizer and the denoiser. Only one image
token per frame is drawn, the model uses $N$. The resolution in (a) is the one
the world model was trained at. The probe in this section crops to
$256\times256$ instead, which gives a $16\times16$ grid.}
\label{fig:navsim-sketch}
\end{figure}

\subsection{Two independent contexts}
\label{sec:navsim-contexts}

The model carries \emph{two} unrelated notions of context, and they are easy to
conflate even though the code keeps them separate. The \emph{trajectory context}
$\texttt{context\_traj}=c$ is the number of leading frames whose ground-truth
velocities are handed to the model and therefore not predicted. Those frames are
copied into the noisy trajectory before the network runs, held fixed during ODE
integration, and left out of the loss. The \emph{image context}
$\texttt{context\_images}=T_{\text{img}}$ is the number of leading frames that
carry an image latent. The two are set independently and the episode length $T$
is free, so everything after the trajectory context gets predicted no matter how
long the episode is.

The runs reported here use $c=0$ and $T_{\text{img}}=2$. This is the strictest
setting for a probe, and the most informative one: the model gets two
front-camera observations and \emph{no} ground-truth motion at all, and has to
produce the whole trajectory from the frozen latent. It also keeps the question
from Section~\ref{sec:probe} free of trajectory extrapolation. With $c>0$ a
constant-velocity continuation of the given prefix would already score well, so
$c=0$ is what isolates the image representation. The context branch is still kept
alive in the architecture, because the planned experiments include training with
the context trajectory supplied half the time. With $c=0$ the context MLP just
receives an empty slice and contributes no tokens.

One implementation detail looks like information leakage on a first reading and
is not: the sampler passes the ground-truth velocities to the denoiser at
inference time. They are only used to fill the first $c$ entries of the
trajectory, so with $c=0$ they are never read. The interface is written this way
so that one call site serves both settings.

\subsection{Online velocity normalization}
\label{sec:navsim-norm}

Section~\ref{sec:syn-norm} showed why per-channel velocity standardization
matters: an isotropic diffusion prior does not match two channels of very
different scale, and the mismatch turns into endpoint drift after the cumulative
sum. On real logs the argument is if anything stronger, since the longitudinal
velocity of a car at $2$\,Hz is on the order of metres per frame while the
lateral component is near zero on straights. But the one-shot computation used in
the synthetic study is inconvenient here, because the statistics would have to be
recomputed every time the dataset, the split or the frame rate changed.

The production path therefore keeps the same transform and estimates its
parameters online. A small module maintains an exponential moving average of the
per-channel mean and variance with decay $0.999$, seeded directly by the first
batch it sees,
\begin{equation}
\mu \leftarrow \gamma\,\mu + (1-\gamma)\,\mu_{\text{batch}},
\qquad
\sigma^2 \leftarrow \gamma\,\sigma^2 + (1-\gamma)\,\sigma^2_{\text{batch}},
\qquad \gamma = 0.999,
\label{eq:ema}
\end{equation}
and applies the familiar pair
$\hat v=(v-\mu)/(\sigma+\varepsilon)$ and $v=\hat v(\sigma+\varepsilon)+\mu$.
Three details make this work here. The batch statistics are computed over the
\emph{valid} frames only, using the same mask as the loss. Without that, the
zero-padded tail of every short episode would drag the mean and the variance
towards zero, and the more ragged the batch the worse the distortion. The
statistics are stored as non-trainable parameters, so they are part of the
checkpoint and a reloaded model denormalizes the same way the trained one did.
That is easy to get wrong when the statistics live in the dataset instead. And
the update only happens while the module is in training mode, so validation and
sampling read the statistics without disturbing them.

\subsection{Training, sampling and configuration}
\label{sec:navsim-training}

The objective is the same as in Section~\ref{sec:syn-flow}: a linear
flow-matching schedule with $\alpha(t)=1-t$ and
$\sigma(t)=\sigma_{\min}+t(1-\sigma_{\min})$, one diffusion time drawn uniformly
per episode, and a mean-squared regression onto the velocity field
$u=x-(1-\sigma_{\min})\varepsilon$. Two things differ. The average is the masked
one of Equation~\eqref{eq:masked-loss}, and only the frames past the trajectory
context contribute.

Sampling mirrors the synthetic procedure with the context handling added. The
trajectory is initialized from pure noise, the first $c$ entries are overwritten
by the known context velocities, and the ODE is integrated backwards from $t=1$
to $t=0$ with $50$ uniform Euler steps, updating only the predicted suffix at
each step. The integrated velocities are denormalized with the EMA statistics and
turned into positions by a cumulative sum re-anchored at the first frame. By
Equation~\eqref{eq:navsim-velocity} that places the predicted path at the origin
in the ego frame of the first observed frame, which is the same frame the cached
trajectory and the map BEV live in, so predictions, ground truth and map overlay
with no further transform.

The validation step logs three quantities of different character on purpose.
\texttt{val/loss} is the \emph{same} masked flow-matching objective as
\texttt{train/loss} evaluated on held-out episodes, so the two curves can be read
against each other and the generalization gap is visible at a glance. It costs
one forward pass. \texttt{val/mse} and \texttt{val/std} need full ODE rollouts
instead, so they measure the sampled trajectories rather than the regression
target. Table~\ref{tab:navsim-config} lists the configuration.

\begin{table}[htbp]
\centering
\caption{Configuration of the NAVSIM probe. The values are those of
\texttt{exp\_navsim/config.yaml} for the runs discussed here. The episode length
$T$ is not a hyperparameter; it is a property of each batch.}
\label{tab:navsim-config}
\begin{tblr}{
  colspec={llX[l]},
  row{1}={font=\bfseries},
}
\toprule
Parameter & Value & Role \\
\midrule
source dataset        & navtrain     & full NAVSIM training logs, $\sim\!40$-frame scenes \\
frame rate            & $2$\,Hz      & stored rate; no subsampling \\
episodes (cached)     & $9\,577$     & longest on-disk contiguous run per scene \\
image size            & $256\times256$ & resize $+$ centre crop of \texttt{CAM\_F0} \\
latent stream         & \texttt{encoded\_q\_sem} & DINOv2-regularized semantic latents \\
latent channels $C$   & $16$         & tokenizer output width \\
\midrule
$\texttt{context\_traj}$   & $0$     & no ground-truth motion given; all frames predicted \\
$\texttt{context\_images}$ & $2$     & observed frames carrying a latent \\
$\texttt{image\_tokens}$   & $9$     & tokens the spatial block emits per observed frame \\
hidden width $D$      & $64$         & embeddings, attention and feed-forward width \\
temporal blocks       & $2$          & cross-attention denoiser stack \\
attention heads       & $8$          & per block, dropout $0.1$ \\
\midrule
$\sigma_{\min}$       & $10^{-5}$    & flow-matching schedule floor \\
diffusion steps       & $50$         & Euler steps at sampling time \\
EMA decay             & $0.999$      & online velocity statistics \\
batch size            & $100$        & episodes per step, zero-padded to $T_{\max}$ \\
base LR               & $5\cdot10^{-4}$ & AdamW, weight decay $0.01$, cosine decay to $0.1\times$ \\
warmup                & $0$ steps    & \\
precision             & 16-mixed     & \\
validation samples    & $20$         & rollouts per context, for \texttt{val/std} \\
validation interval   & every $5$ epochs & \\
\bottomrule
\end{tblr}
\end{table}

\subsection{Metrics for a distribution of futures}
\label{sec:navsim-metrics}

A probe that is asked for a \emph{distribution} cannot be judged by pointwise
error alone, because the two failure modes we care about sit at opposite ends of
the same axis. A model that collapses onto the mean of the data will look
respectable under a squared error and be useless as a world-model read-out. A
model that samples wildly will cover every future and be just as useless. So we
log two quantities.

The first is the masked trajectory MSE, averaged over the predicted and valid
frames of all sampled rollouts, which says how close the samples sit to the one
future that actually happened. The second measures spread. The \emph{total turn}
of a trajectory is the accumulated change of heading along it: with segments
$s_t=p_{t+1}-p_t$ and headings $\theta_t=\operatorname{atan2}(s_{t,y},s_{t,x})$,
\begin{equation}
\Theta \;=\; \sum_{t} \operatorname{wrap}\!\big(\theta_{t+1}-\theta_t\big),
\qquad
\operatorname{wrap}(\delta)=\operatorname{atan2}(\sin\delta,\cos\delta),
\label{eq:total-turn}
\end{equation}
The wrapping keeps the increments in $[-\pi,\pi]$, so a heading crossing the
$\pm\pi$ branch cut does not register as a full turn. A turn increment at index
$t$ is built from two consecutive segments and therefore from three consecutive
frames, so an increment is only counted when the third of those frames is valid.
The padded tail contributes nothing, which is what keeps the metric comparable
across episodes of different length. Sampling the same context $N=20$ times,
taking the standard deviation of $\Theta$ across those samples, and averaging
over the batch,
\begin{equation}
\texttt{val/std} \;=\; \frac{1}{B}\sum_{b} \operatorname{std}_{n\le N}\ \Theta_{b}^{(n)},
\label{eq:turn-std}
\end{equation}
gives an interpretable read-out of how much the probe disagrees with itself about
where the car is going. The pair has to be read together. A low MSE with a
near-zero turn spread is the signature of mode collapse onto the average path,
while a growing spread that does not hurt the MSE is evidence of real
multi-modality.

The qualitative counterpart is a per-episode PDF, produced by a separate
evaluation script from the same deterministic validation split. Each page shows
the first and last context front views, the NAVSIM map bird's-eye view with the
\emph{whole} sampled distribution drawn on it, and the trajectory panels. The
model predicts the full padded length for every sample, so the page trims each
episode back to its true length before plotting. Episodes are pre-filtered by
driven distance ($\ge 5$\,m) and by the angle between the initial heading and the
start-to-end vector ($\ge 10^\circ$). A stationary or perfectly straight episode
makes any trajectory distribution look correct, so the filter is there to pick
out the cases where the question is not trivial.

\subsection{Status, limitations and planned experiments}
\label{sec:navsim-status}

The pipeline above is the state the system currently runs in: whole-length
episodes, padded batches with a masked objective, the multi-token image
aggregator feeding a cross-attention denoiser, online velocity normalization, and
a validation protocol reporting both accuracy and spread. Going from fixed short
windows to full episodes, together with the aggregator revision, is what moved
the qualitative predictions from unconvincing to clearly structured.

Several limitations are known and should be stated plainly. The attention stack
does not yet get a key-padding mask, so padded trajectory tokens are visible as
keys even though the loss excludes them. Those tokens are zero-valued and the
effect looks small, but it is a real inconsistency between the loss and the
attention, and it is the next thing to fix. The image context is short
($T_{\text{img}}=2$) and its frame identity enters only through AdaLN modulation,
as discussed in Section~\ref{sec:navsim-aggregator}. The probe is evaluated with
MSE and turn spread rather than with the displacement-based metrics the driving
literature uses, and there is still no control baseline, meaning a
constant-velocity extrapolation or a probe trained on a shuffled latent, to
establish how much of the performance comes from the representation rather than
from the strong prior structure of driving trajectories. Finally, and this
matters most for the claim the project set out to test, everything reported here
conditions on \emph{stage-1 tokenizer} latents. The stage-2 world model of
Section~\ref{sec:orbis} has not been probed at all.

That last gap is what the planned continuation is about, and it has three parts.
The first is a layer-wise probe of the world model: freeze the STDiT stack, tap
its intermediate block activations, and train the same probe on each tap in turn
to find where ego-motion information is strongest as a function of depth. A flat
depth curve and a peaked one would both tell us something. The second is a
capacity and control sweep that fixes the best tap and varies probe size against
baselines, so that a good score can be credited to the representation rather than
to the read-out. The third goes at calibration directly: sample
$k\in\{1,5,10,20\}$ futures per context on strata separating free-flow driving
from branching intersections, and report $\text{minADE}_k$ and mode recall on map
branches next to a sharpness measure on straights, so a model cannot score well
just by sampling everything. That third experiment needs longer and denser source
clips, since a five-second NAVSIM scene rarely contains a real branch point, and
the data ingest for that is the one piece of infrastructure still missing.
